%! AP Statistics — Unit 2, Lesson 2.1 — Student Packet Cover Sheet
\documentclass[10pt]{article}
\usepackage{apstats-article}
\usepackage{apstats-boxes}

%=============================================================================
\begin{document}

% ── Full-bleed navy banner ────────────────────────────────────────────────────
\begin{tikzpicture}[remember picture, overlay]
  \fill[navy] (current page.north west)
    rectangle ([yshift=-0.9in]current page.north east);
\end{tikzpicture}
\vspace{-0.8in}

\begin{center}
  {\color{white}\textbf{\LARGE AP Statistics}}\\[4pt]
  {\color{skymid}\large\textbf{Unit 2: Exploring Two-Variable Data}}\\[6pt]
  {\color{white}\textbf{\large Lesson 2.1 \quad Introducing Statistics: Are Variables Related?}}
\end{center}

\vspace{0.22in}

\namedateperiod

\vspace{0.18in}

% ── LEARNING TARGETS ─────────────────────────────────────────────────────────
\begin{learningtargetbox}
\small
\begin{enumerate}[leftmargin=1.5em, itemsep=5pt]
  \item \ldots{} identify a statistical question that asks about a possible relationship
        between two variables.
  \item \ldots{} explain why an apparent pattern or association in data may be random and
        not a real relationship.
  \item \ldots{} explain why association does not imply causation, and give an example of
        a confounding variable.
\end{enumerate}
\end{learningtargetbox}

\vspace{0.14in}

% ── TABLE OF CONTENTS ────────────────────────────────────────────────────────
\begin{tocbox}
\small
\renewcommand{\arraystretch}{1.5}
\begin{tabularx}{\linewidth}{@{} c l X r @{}}
\rowcolor{sky}
\textbf{\#} & \textbf{Component} & \textbf{Description} & \textbf{Score} \\
\midrule
1 & Warm-Up       & Spiral review + two-variable thinking   & \blank{1.2cm} \\
2 & Guided Notes  & Statistical questions, association, causation & \blank{1.2cm} \\
3 & Group Activity & Evaluating real vs.\ random associations & \blank{1.2cm} \\
4 & Exit Ticket   & Analyze a study for association and confounding & \blank{1.2cm} \\
5 & Homework      & News blurbs + association analysis        & \blank{1.2cm} \\
\midrule
  & \multicolumn{2}{r}{\textbf{Total}} & \blank{1.2cm} \\
\end{tabularx}
\end{tocbox}

\vspace{0.14in}

% ── KEEP IN MIND ─────────────────────────────────────────────────────────────
\begin{remindbox}
\small
These are the big ideas of Unit 2 --- and of AP Statistics. Start building them now.

\vspace{6pt}
\begin{tabularx}{\linewidth}{@{} >{\centering\arraybackslash\columncolor{goldbg}}p{0.06\linewidth}
                                    X @{}}
\textbf{\large!} &
\textbf{Apparent patterns may be random.}\enspace
A pattern visible in a small sample might disappear if you collected more data.
Statistics gives us tools to judge whether a pattern is real. \\[6pt]
\textbf{\large!} &
\textbf{Association $\neq$ Causation.}\enspace
Two variables can move together without one causing the other.
A third variable --- a \emph{confounding variable} --- may explain both. \\[6pt]
\textbf{\large!} &
\textbf{The question determines the method.}\enspace
Two categorical variables call for a different display than two quantitative variables.
Choosing the right method starts with identifying the right question. \\[6pt]
\textbf{\large!} &
\textbf{Unit 2 builds on Unit 1.}\enspace
Everything you learned about describing one variable now applies to each variable in a pair.
The new question is: \emph{do they move together?} \\
\end{tabularx}
\end{remindbox}

\end{document}
